Het brak haar laatste wil.
“Ik blijf voor altijd stil!”,

was haar laatste protest.
“Dat zal worden getest!
Je zal worden gepest!

Ik zal je buik vullen.
Stop dus met het lullen
over flutte prullen.

Het was toch jouw besluit?
Dus lever ik je uit
of blijf je toch mijn buit?”

Ze kroop tegen hem aan.
Hij kon haar niet verstaan.
Haar mond liet geen woord gaan.

Verstomd door haar trauma
verstopt in haar vulva.
Ze zweeg tegen opa.

Hij was nota bene,
waarschijnlijk degene,
die haar psychogene

mutisme had bezorgd.
Wie dan leeft, die dan zorgt.
Zijn afspraak gewaarborgd;

“In mijn mooie thuisstad
breng ik jou als kerkschat.
Daar geldt de regel dat

je gevrijwaard bent
en compleet wordt erkend,
zolang een vaste vent

zich over je ontfermt.
Al ben je uitgezwermd,
je blijft bij ons beschermd.

Je bestaan wordt wel hels.
Je hapte niet rebels,
dus hier heb je een pels.”

Ze kreeg een schapenvacht,
die hij altijd meebracht
om, als hij overnacht,

droog op te gaan slapen.
Emmy hield van schapen,
maar moest toch dit wapen

tegen de kou innen.
Die drong te sterk binnen.
Zo wist hij te winnen.

Helaas was de schaapsvacht,
die de jager inbracht,
te kort om haar lijfwacht

geen zicht te verlenen
op haar pruim en benen.
Die hij vrij kon lenen

met zijn grijpgrage hand
om een driftige brand,
die in hem was beland,

door naar haar te turen,
nog wat aan te vuren.
Zo verliepen uren.

Zijn bruinbierconsumptie en dehydratatie verwonderde Emmy.

Het verhaal van Psyche en haar zoektocht naar Eros is een archetypisch verhaal dat diep geworteld is in de menselijke ervaring van verlangen, liefde, en innerlijke transformatie. In het kader van het moeilijke takenmotief wordt de reis van Psyche naar haar geliefde een symbolische weergave van de uitdagingen en beproevingen die we moeten doorstaan om innerlijke eenheid en harmonie te bereiken.

Psyche, een sterfelijke vrouw van ongeëvenaarde schoonheid, wordt door de god van de liefde, Eros, begeerd. Eros neemt haar mee naar een prachtig paleis waar ze in het donker samen zijn, zodat Psyche zijn goddelijke gedaante niet kan zien. Wanneer ze zich overgeeft aan de nieuwsgierigheid, steekt ze een lamp aan en ziet Eros in al zijn goddelijke schoonheid. Maar een druppel olie van de lamp valt op hem, en Eros ontwaakt en vlucht, teleurgesteld door haar wantrouwen.

Om Eros terug te winnen, moet Psyche vier onmogelijke taken volbrengen, opgelegd door Aphrodite, de jaloerse moeder van Eros. Deze taken staan symbool voor de innerlijke reis die ze moet maken om volwassen te worden en haar liefde waardig te zijn.
Het Sorteren van de Zaden
De eerste taak is het sorteren van een enorme berg zaden – gerst, tarwe, maanzaad, enzovoort – binnen één nacht. Deze taak lijkt onoverkomelijk, maar ze krijgt hulp van een leger mieren die haar helpen. Dit symboliseert de noodzaak om onderscheid te maken en orde te scheppen in de chaos van het leven. Het vraagt om een nauwkeurig onderscheidingsvermogen, een taak die vaak onmogelijk lijkt zonder hulp van onbewuste krachten (zoals de mieren).
De Wol van het Gouden Schaap
De tweede taak is het verzamelen van wol van gevaarlijke, gouden schapen. Psyche leert echter dat ze deze wol moet verzamelen van takken waar de schapen langs gewreven hebben, zonder zichzelf in gevaar te brengen. Dit gaat over het leren kennen van de juiste timing en aanpak; niet alles moet frontaal aangegaan worden. Soms moet je wachten op het juiste moment, en op een slimme manier handelen.
Water uit de Rivier Styx
De derde taak is het verzamelen van water uit de dodelijke rivier Styx. Deze taak symboliseert de confrontatie met de dood en het onbekende. Psyche krijgt hulp van een adelaar, een symbool voor een hoger perspectief. Dit wijst op de noodzaak om een bredere visie te ontwikkelen en boven het alledaagse uit te stijgen om de angsten van het onbewuste onder ogen te zien.
De Schoonheid van Persephone
De laatste taak is het verkrijgen van een doosje met de schoonheid van Persephone uit de onderwereld. Deze taak vraagt om een afdaling in de diepste lagen van het zelf, een confrontatie met de donkerste angsten. Uiteindelijk kan Psyche de verleiding niet weerstaan om het doosje te openen, wat leidt tot een diepe slaap (vergelijkbaar met de dood). Hier wordt duidelijk dat zelfs de sterkste wilskracht faalt zonder nederigheid en overgave.
Het moeilijke takenmotief in dit verhaal weerspiegelt de innerlijke beproevingen die we allemaal moeten doorstaan op weg naar zelfontplooiing. Elke taak symboliseert een stap in de evolutie van Psyche's ziel, van onschuld en onwetendheid naar wijsheid en liefde. De taken die eerst onoverkomelijk lijken, worden hanteerbaar door hulp van onverwachte bondgenoten en het ontwikkelen van inzicht. Dit laat zien dat groei niet mogelijk is zonder uitdagingen en dat we soms moeten vertrouwen op krachten buiten onszelf.

De reis van Psyche is er een van transformatie. Ze begint als een meisje dat wordt gedreven door verlangen en nieuwsgierigheid, maar leert door haar beproevingen wat liefde werkelijk betekent. Ze moet haar angsten en zwakheden onder ogen zien, leren loslaten, en uiteindelijk zichzelf overgeven aan iets dat groter is dan zijzelf. Het is deze reis die haar waardig maakt om herenigd te worden met Eros, niet slechts als geliefde, maar als gelijke. Hun vereniging aan het eind symboliseert de hereniging van het mannelijke en vrouwelijke, het bewuste en onbewuste, de psyche en de geest.
Het verhaal leert ons dat echte liefde en innerlijke heelheid pas bereikt kunnen worden na het doorstaan van beproevingen die ons uitdagen om te groeien en te transformeren. Het moeilijke takenmotief benadrukt dat deze uitdagingen niet slechts obstakels zijn, maar noodzakelijke stappen op de weg naar ons diepste zelf.

et woord “psyche” zoals wij dat in psychologische termen gebruiken, is inderdaad afgeleid van het Griekse woord ψυχή (psuchē), wat oorspronkelijk “ziel”, “geest” of “levensadem” betekende. In de klassieke Griekse cultuur verwees het naar het onsterfelijke deel van een persoon, dat hun gedachten, emoties en identiteit omvatte.
Het verhaal van Psyche in de mythologie is een symbolische representatie van deze betekenis. Psyche personifieert de ziel van de mens, die verlangt naar eenheid en harmonie met de goddelijke liefde, gesymboliseerd door Eros. Haar beproevingen en uiteindelijke vereniging met Eros weerspiegelen de reis van de menselijke ziel door uitdagingen en groei naar verlichting en vervulling.
In de moderne psychologie wordt het begrip “psyche” gebruikt om het geheel van iemands mentale en emotionele processen te beschrijven. Het omvat alles wat te maken heeft met de geest, zoals gedachten, emoties, herinneringen en verlangens. Zo wordt bijvoorbeeld de term “psychologie” afgeleid van “psyche” en “logos” (woord, rede of studie), wat letterlijk de “studie van de ziel” betekent.
De overgang van de mythologische betekenis naar de wetenschappelijke term weerspiegelt een verschuiving van een spirituele naar een meer wetenschappelijke benadering van de menselijke geest, maar de kern – het begrijpen van de diepere aard van menselijk bewustzijn en gedrag – blijft vergelijkbaar.
